\chapter{План на двенадцать недель}

\begin{warnbox}[Глава заменена роадмапом экосистемы]
Состояние на 29 сентября 2026 года. План на двенадцать недель не исполняется. Порядок работ задаёт список приоритетов роадмапа экосистемы (вики, \texttt{roadmap/ecosystem-roadmap.ru.md}, черновик v4, §2b): первым идёт модуль хоста сессий и на нём waggle CM1, вторым конфигурация процесса PC0--PC1 с первыми демонстрационными плейбуками, дальше этапы M0--M8 и треки CM, PC, KP и AR. Сроки и готовность каждого шага ведут треки в репозиториях модулей. Текст главы сохранён как история; где он расходится с роадмапом, прав роадмап.
\end{warnbox}

Цель: через три месяца работающий MVP, который соединяет существующую инфраструктуру в продукт и готов к обкатке на реальных проектах. Пилот не есть отдельный этап: метрики встроены в MVP, и обкатка начинается в день готовности. Уровень доверия первого развёртывания L0, с определением кворума как параметром, чтобы L1 и L2 не требовали переписывания.

\section{Три блока}

\begin{center}
\begin{ganttchart}[
  hgrid, vgrid={*{3}{draw=none},dotted},
  x unit=2.6mm, y unit chart=6.5mm, y unit title=6.5mm,
  title label font=\footnotesize\sffamily, bar label font=\footnotesize,
  bar/.style={fill=paleblue,draw=accent}, bar height=0.55,
  group/.style={fill=accent,draw=accent}, group label font=\footnotesize\sffamily\bfseries,
  milestone/.style={fill=accentamber,draw=accentamber}, milestone label font=\footnotesize
]{1}{36}
\gantttitle{Недели 1--4}{12}\gantttitle{Недели 5--8}{12}\gantttitle{Недели 9--12}{12}\\
\ganttgroup{Блок 1: исполнение}{1}{12}\\
\ganttbar{Демон среды, macOS и Linux}{1}{10}\\
\ganttbar{Скилл shell, восстановленная сессия}{3}{10}\\
\ganttbar{Минимальный runtime, подпись}{5}{11}\\
\ganttmilestone{Агент проходит полный цикл}{12}\\
\ganttgroup{Блок 2: журнал}{13}{24}\\
\ganttbar{Секвенсер поверх реестра}{13}{21}\\
\ganttbar{Светлый клиент, откат, включение}{16}{22}\\
\ganttbar{Кандидат, patch-id, CAS}{17}{23}\\
\ganttmilestone{Интеграция только по QC}{24}\\
\ganttgroup{Блок 3: рой}{25}{36}\\
\ganttbar{Аренда, планировщик, области}{25}{33}\\
\ganttbar{Случайный ревьювер, три оси}{27}{33}\\
\ganttbar{Остановка, join, метрики}{29}{35}\\
\ganttmilestone{MVP, обкатка на проектах}{36}
\end{ganttchart}
\end{center}

\subsection{Блок 1, недели 1--4: исполнение}

Демон среды как провайдер исполнения на хосте: Seatbelt и bubblewrap, санбокс на задачу с клоном дерева, защищённые пути только на чтение, receipts на команду, инвентарь, планировщик слотов, egress-прокси, CA среды с сертификатами на задачу и принудительной командой поверх OpenSSH (встроенный SSH-сервер в демоне откладывается до публичного роя). Скилл shell с восстановленной сессией, именованными сессиями, фоном, spill и переподключением. Минимальный runtime с планом из двух привязок и подписью на машине владельца.

Критерий приёмки: один агент проходит полный цикл на реальном репозитории проекта через среду, а не локально: клонирует, правит, собирает, тестирует, представляет подписанный коммит; транскрипт структурно неотличим от локального прогона; ни одной команды вне санбокса; receipts на каждую команду.

\subsection{Блок 2, недели 5--8: журнал}

Секвенсер поверх реестра аудита: типы записей, включая версии контракта и зависимости, кворум как порог эндорсментов, снапшоты, доказательства включения, сертификаты доступности, сертификат остановки. Среда как светлый клиент: открытие санбокса только по назначению с QC, защита от отката, проверка актуальности на границе исполнения. Интеграционный кандидат до финального одобрения, атомарное продвижение головы, операции раннера с ключами идемпотентности и тремя исходами. Борд как содержимое журнала: переходы как коммиты борда с QC.

Критерий приёмки: ни один санбокс не открывается и ни одна интеграция не проходит без записи с QC; одобрение и аттестации, адресованные не к интегрируемому кандидату или к прежней версии контракта, отвергаются; попытка интегрировать при ушедшей вперёд голове приводит к пересборке кандидата с повторным ревью; повторная подача старого снапшота отвергается; операция с исходом не установлен не повторяется без сверки. Все трассы исполняемой модели проходят как векторы.

\subsection{Блок 3, недели 9--12: рой}

Аренда оркестратора с fencing и heartbeat поверх спавна борда; планировщик по ёмкости, окнам квот, слотам и областям; заморозка кандидатов, случайный ревьювер, три оси; меню при превышении $N$; сертификат остановки и возобновление; команда присоединения по манифесту; дашборд метрик.

Критерий приёмки: три и больше контрибьюторов с ёмкостью разного происхождения одновременно ведут задачи одной цели; ревьювер никогда не совпадает с продьюсером по трём осям; остановка одной подписью прекращает открытие санбоксов и интеграции; смена оркестратора по таймауту не теряет состояния; пять метрик показываются из журнала.

\section{Критический путь}

Демон среды и скилл shell. Всё остальное можно временно заменить заглушками: журнал секвенсером на одной ноде, оркестратора человеком, ревьювера вторым агентом на той же машине. Но без демона со санбоксом и без shell, неотличимого от локального, нет ничего, что можно показать. Поэтому первые две недели уходят только туда, и туда же направляется большинство агентов.

\section{Что отложено без потери формы}

Уровни L1 и L2: кворум есть параметр, и реализация L0 не закрывает путь. Переносимая репутация между проектами: аттестации пишутся с первого дня, взвешивание позже. Экономика и учёт вклада: данные пишутся (оценка, исход, вестинг), правила раздела пулов считаются потом. Полировка передачи сессий между агентами: плановая передача через ax уже есть, тонкая настройка позже. Верификатор как отдельная машина: сначала одна среда, вторая при первой возможности. Реестр федерации: сначала манифест по адресу, реестр когда проектов больше одного. Слой идентичности: типы записей и поле доказательств в join закладываются сейчас, резолверы и подтверждение кворумом валидаторов в первом релизе после MVP.

\section{Как двигаться с агентами}

Скорость достигается тем же циклом, которым проект уже живёт: спецификация впереди реализации, с нормативными формулировками и векторами соответствия, затем агенты реализуют против векторов параллельно. Порядок:

\begin{enumerate}
\item Неделя 1: нормативные документы для демона и для типов записей журнала. Первый набор векторов соответствия получается запуском исполняемой модели фрагмента интеграции, а не ручным написанием; вручную дописываются векторы для санбокса и внешних эффектов. Это единственная неделя, где узкое место люди.
\item Со второй недели: агенты реализуют бэкенды, скилл, секвенсер и клиент против векторов, каждая задача через борд с ревью.
\item Люди держат два места, где нужна голова: модель полномочий (что считать областью, как компилировать план из области) и семантику интеграционного кандидата (что делать при неравенстве patch-id, как строить очередь).
\end{enumerate}

Когда демон и скилл готовы, рой начинает строить сам себя: следующие задачи проекта идут через протокол, и первые метрики появляются не на чужом проекте, а на своём.

\section{Риски}

\begin{center}
\small
\begin{tabular}{p{4.6cm} p{8.8cm}}
\toprule
Риск & Что делаем \\
\midrule
условия вендоров по подпискам в стороннем runtime & ёмкость любого происхождения: ключи API и локальные модели как равноправные источники; подписки там, где допустимо; проверка условий до публичного запуска \\
ёмкость одной среды & слоты и очередь с первого дня; вторая машина как реплика и верификатор при первой возможности \\
ошибки в санбоксе на macOS & инвентарь честно объявляет частичное; никаких общих записываемых каталогов; всё ценное за операциями раннера; тест на побег как часть векторов \\
стимулы ревью (отклонять на всякий случай) & измеряем долю пересмотренных отклонений; правило оплаты раундов и споров настраивается по данным \\
шторм перебазирований и повторных ревью & дизъюнктные области как правило валидатора; декомпозиция по областям; стоимость повторного ревью кандидата измеряется, перенос одобрения возможен только как расширение с условием по прочитанным путям \\
терминология (три значения слова эпоха) & глоссарий в спецификациях; в коде разные имена: plan epoch, lease epoch, orchestrator epoch \\
\bottomrule
\end{tabular}
\end{center}
